\section{The halving graphs}
\label{app:halving}

These graphs show that Theorem~U is sharp and supply the literal examples
in Corollary~L. Let $f=2^t$ with $t\ge0$. For a positive integer~$i$ let
$\rho(i)$ be the exponent of the largest power of~2 that divides~$i$ (the
``ruler function''; for example $\rho(12)=2$). The graph has girls~$u_i,c_i$
and boys~$b_i,d_i$ for $1\le i\le f$, and the edges
\[
  u_ib_i,\quad c_ib_i,\quad c_id_i\quad(1\le i\le f),\qquad
  u_i\,b_{i+2^{\rho(i)}}\quad(1\le i<f).
\]
So $M=N=s=2f$ and $E=4f-1$; the edges $u_ib_i$ and $c_id_i$ form a perfect
matching.

\paragraph{The input.} Number $b_i$ as $2i-1$ and $d_i$ as $2i$, and number
$c_i$ as $f+i$, with list $[b_i,d_i]$. Sort the indices~$i$ by decreasing
$\rho(i)$, breaking ties by increasing~$i$, and give the corresponding
girls~$u_i$ the numbers $f,f-1,\ldots,1$ in that order. The list of $u_i$ is
$[b_i,b_{i+2^{\rho(i)}}]$, or $[b_f]$ when $i=f$. For the three literal
examples of Corollary~L this gives
\begin{center}\small
\begin{tabular}{cl}
\toprule
$s$&girl lists, in increasing order of girls\\
\midrule
2&$[1];\ [1,2]$\\
4&$[1,3];\ [3];\ [1,2];\ [3,4]$\\
8&$[5,7];\ [1,3];\ [3,7];\ [7];\ [1,2];\ [3,4];\ [5,6];\ [7,8]$\\
\bottomrule
\end{tabular}
\end{center}

\paragraph{The run.} Step~\Hs1 couples every $c_i$ with $b_i$ and leaves all
$u_i$ free, in the order just described. After $r$ dag rounds
($0\le r\le t$) the free girls are the $u_i$ with $2^r\mid i$, and the
indices fall into \emph{blocks} of $2^r$ consecutive integers
$i-2^r<j\le i$, one block for each free girl~$u_i$. Her basic path visits
$r+1$ of the girls~$u_j$ of her block, namely $u_i$ and then $u_j$ for
$j=i-2^r+2^q$ with $q=r-1,\ldots,1,0$, and ends with $c\sg d$ at the
smallest index of the block; it has level~$r+2$. For example, when $f=4$
and $r=1$ the free girls are $u_2$ and $u_4$, and the basic path of $u_4$
is $u_4\sg b_4\db u_3\sg b_3\db c_3\sg d_3$.

Every augmenting path has level~$r+2$ at this stage: a free girl can start
along $b_i$ or along $b_{i+2^{\rho(i)}}$, and after that first step the
alternating continuation is forced. In the next round step~\Hs6 lists, at
each $b_i$, the free girls in reverse queue order. As a result the search
joins each block whose free girl~$u_i$ has $\rho(i)=r$ to the block on its
right; the other roots fail at deleted girls; and step~\Hs{16} removes
exactly the girls with $\rho(i)=r$, which form the end of the queue. So the
free girls go $f\to f/2\to\cdots\to1\to0$ in $t+1$ dag rounds, and
property~$\Up$ holds in every round. (For $f\le32$ we have also checked
this by simulation, including a census of all augmenting paths.)

\paragraph{Work.} The round that starts with $n'=f/2^r$ free girls has
$W_5=2n'(r+2)-1$ and $W_{13}=W_{11}=n'(r+2)$; in total
$\sum W_5=12f-3t-9$ and $\sum W_{13}=\sum W_{11}=6f-t-4$, and~\Hs1 examines
$3f-1$ entries. So these graphs need only linear work in all.
